\documentclass[12pt]{article}
\usepackage[margin=0.72in]{geometry}
\usepackage{lmodern}
\usepackage{microtype}
\usepackage{multicol}
\usepackage{fancyhdr}
\usepackage{titlesec}
\usepackage{enumitem}
\usepackage{xcolor}
\usepackage[hidelinks]{hyperref}

\setlength{\columnsep}{0.30in}
\setlength{\parindent}{1.05em}
\setlength{\parskip}{0.24em}
\setlength{\headheight}{15pt}
\titleformat{\section}{\large\bfseries\scshape}{}{0pt}{}
\titleformat{\subsection}{\normalsize\bfseries}{}{0pt}{}
\pagestyle{fancy}
\fancyhf{}
\fancyhead[L]{Whately: Terms and Meaning}
\fancyhead[R]{Reading Handout}
\fancyfoot[C]{\thepage}
\renewcommand{\headrulewidth}{0.4pt}

\newcommand{\focusbox}[1]{%
\vspace{0.4em}\noindent\colorbox{gray!12}{\begin{minipage}{0.96\linewidth}#1\end{minipage}}\vspace{0.5em}}

\begin{document}
\begin{center}
{\LARGE\bfseries Logic Is Not Magic}\par\vspace{0.15em}
{\large\scshape Reading 2: One Word, Two Ideas}\par\vspace{0.45em}
\rule{0.82\textwidth}{0.4pt}\par\vspace{0.6em}
\end{center}

\section*{Vocabulary Preview}
\begin{multicols}{2}
\begin{itemize}[leftmargin=*, itemsep=0.18em]
\item \textbf{Term}: A word or phrase used to stand for a particular idea.
\item \textbf{Meaning}: The specific idea a term refers to in a given context.
\item \textbf{Ambiguity}: The condition in which one word or phrase can refer to more than one idea.
\item \textbf{Equivocation}: Using one word in two different senses within a single argument.
\item \textbf{Fix the term}: To state clearly which meaning is intended before evaluating an argument.
\item \textbf{Context}: The surrounding words and situation that help determine meaning.
\end{itemize}
\end{multicols}

\section*{Introduction}
In the first reading we learned to slow an argument down and ask four questions: What is the conclusion? What reason is actually given? What hidden assumption must be true? Does the conclusion really follow?

One of the most common reasons arguments fail this test is that a key word is being used in two different senses without anyone noticing the shift. Richard Whately observed that many faulty conclusions and pointless disputes arise not from bad reasoning, but from words that quietly change their meaning in the middle of an argument.

This reading is adapted from Whately's treatment of terms and ambiguity in \textit{Elements of Logic}. The wording is modernized, but the central concern remains Whately's: reasoning becomes unstable when words are allowed to carry more than one idea inside the same argument.

\focusbox{\textbf{Source note:} Adapted and modernized from Richard Whately's treatment of terms and ambiguity in \textit{Elements of Logic}. This version preserves Whately's central claim that many apparent failures of logic result from unacknowledged shifts in the meaning of terms.}

\begin{multicols}{2}
\section*{Reading}

\subsection*{1. Terms Stand for Ideas}
A \textit{term} is a word or short phrase that stands for an idea. When we say “bank,” we might mean a financial institution or the land beside a river. When we say “light,” we might mean “not heavy” or “bright.” Most words have more than one possible meaning. In ordinary conversation the surrounding words usually make the intended meaning clear.

\subsection*{2. The Problem of Shifting Meaning}
The trouble begins when someone uses the same word in two different senses inside one continuous argument without acknowledging the change. This is called \textit{equivocation}. The argument appears to prove its conclusion only because the meaning of the term has shifted.

\subsection*{3. A Simple Case}
Consider this argument: “No bank is entirely safe from thieves. The edge of a river is a bank. Therefore the edge of a river is not entirely safe from thieves.”

The conclusion is obviously false. The argument collapses once we notice that “bank” is used first in the financial sense and then in the geographical sense. The two meanings were never fixed.

\subsection*{4. A More Serious Example}
A more important case appears in debates over education. Someone argues: “A liberal education does not prepare students for employment. Therefore liberal education is of little value.”

Here the word “liberal” moves between two ideas. In one sense it refers to a broad education in the arts and sciences aimed at intellectual formation. In another sense it refers to a particular political orientation. If the term is fixed to the educational meaning, the claim that such an education lacks value becomes much harder to defend.

\subsection*{5. Equality and Equal Treatment}
The same problem appears in arguments about equality. Someone claims: “We must treat all students equally. Therefore every student’s grade should be equal on every assignment.”

Here the idea of “equal” shifts between two different meanings. In one sense, equal treatment means giving every student the same rules, expectations, and opportunity to demonstrate what they have learned. In another sense, equal treatment is taken to mean producing the same result for every student. The first meaning allows teachers to respond to real differences in preparation and effort. The second meaning makes honest evaluation nearly impossible. The argument only appears fair because the meaning of equal treatment was never fixed.

\subsection*{6. Tolerance and Judgment}
A similar difficulty arises in arguments about tolerance. Someone claims: “We must be tolerant of different beliefs. It is therefore wrong to say that any religion or moral view is mistaken or harmful.”

In this case, the word “tolerant” does not need to be repeated for the meaning to shift. The first sentence uses “tolerant” to mean allowing people to hold and practice their own beliefs without punishment or coercion. The second sentence uses the same idea to mean that one must never judge any belief as false or damaging. These are two different meanings. The argument only appears reasonable because the meaning of tolerance was allowed to change without being noticed.

\subsection*{7. The Practical Point}
Before testing whether a conclusion follows from its reasons, one must first determine whether the key terms have been used in a single, consistent sense throughout the argument. If they have not, the meaning of the disputed term must be fixed before any further evaluation can be fair.

Fixing a term does not require choosing a permanent or perfect definition. It requires only that, for the purposes of the present argument, the speaker states clearly which idea the word is being used to represent. Once the meaning is fixed, the hidden assumption often becomes visible and the four-step test can proceed.

\subsection*{8. Why This Matters}
Readers who learn to detect this kind of ambiguity acquire a defensive habit. They become less easily persuaded by arguments that succeed only because a word has been allowed to mean two things at once. They also become more capable of constructing arguments that remain clear because their central terms have been defined at the outset. This habit will matter throughout our recurring Shakespeare cases. In \textit{Macbeth}, words such as \textit{fate}, \textit{manhood}, and \textit{safety} become dangerous when their meanings slide. In \textit{Julius Caesar}, words such as \textit{honor}, \textit{ambition}, \textit{liberty}, and \textit{tyranny} can carry different meanings for different speakers and audiences.

\section*{Reading Focus}
\begin{enumerate}[leftmargin=*]
\item What is a term, and why can the same word stand for different ideas in different contexts?
\item Using one of the examples in the reading, explain the difference between ambiguity and equivocation.
\item In the example about liberal education, what two distinct ideas does the word “liberal” represent?
\item Why does fixing the meaning of a key term often make the hidden assumption of an argument easier to identify?
\item Choose one word from \textit{Macbeth} or \textit{Julius Caesar}---such as fate, manhood, honor, ambition, liberty, or tyranny. What two meanings could the word carry in an argument?
\end{enumerate}

\section*{Window Note and Summary}
Create a portfolio window note that shows how one argument depends on an unacknowledged shift in the meaning of a key term. You may use an example from this reading, \textit{Macbeth}, or \textit{Julius Caesar}. Include at least two drawings or symbols that represent the two different meanings. At the bottom, write a 5--7 sentence summary explaining why fixing terms matters for evaluating arguments.

\section*{Connection to Week 1}
How does the skill of fixing terms connect to the four questions we practiced in the first reading (conclusion, stated reason, hidden assumption, test)?

\end{multicols}
\end{document}